% Clinical Background section. Paste into your own thesis, or \input it.
%
% FIGURE WANTED: either a CCTA acquisition schematic (tube and detector rotating,
% ECG gating) or a side-by-side of a clean vessel and one degraded by motion and by
% blooming. Not added yet: an external figure needs a licensed credit generated into
% figures/external/CREDITS.md by scripts/fetch_external_figures.py, which does not
% exist on this branch. Do not add a \ref until the figure and its credit exist.

\section{Coronary computed tomography angiography}
\label{sec:ccta}

CCTA is the non-invasive imaging test recommended as the first-line investigation in patients with
suspected chronic coronary syndrome \cite{knuuti2019esc}. The acquisition is brief.
The patient is positioned on the table of a \emph{computed tomography} \textbf{(CT)} scanner, and a
\emph{contrast agent} (a solution containing iodine, which absorbs X-rays strongly) is
injected into a vein in the arm.
The X-ray tube and its detector then rotate around the patient, recording a projection of the chest at each angular position, and those projections are reconstructed into a single three-dimensional volume \cite{abbara2016scct}.
Timing determines whether the scan is usable: acquisition is timed to coincide with the passage of the contrast agent through the coronary arteries, and the result is a volume in which the lumen stands out brightly against the surrounding tissue.

\subsection{Acquisition}

The limits of CCTA as a measuring instrument bound every geometric claim made from these images, and the first of them is motion.
The coronary arteries move with the myocardium throughout the cardiac cycle.
Acquisition is therefore synchronized to the electrocardiogram, and the volume is reconstructed from
the phase of least motion, conventionally mid-diastole \cite{abbara2016scct}.
Heart rate is routinely lowered pharmacologically before the scan, which lengthens that window
\cite{abbara2016scct}.
Where the synchronization fails, or the rhythm is irregular, \emph{motion artifact} blurs the
vessel or duplicates it \cite{abbara2016scct}.


\subsection{Resolution}

Resolution bounds the measurement.
A reconstructed volume is a grid of \emph{voxels}, the three-dimensional equivalent of pixels, and
no structure smaller than a voxel can be resolved.
Proximal coronary arteries measure a few millimeters in diameter and taper distally toward 1 mm and
below.
A vessel spanning only two or three voxels cannot support a reliable radius estimate, which
places a floor on the branch order at which geometric features remain meaningful.
